\begin{samepage}
\bigskip
\pdfbookmark[0]{Magnetische Beobachtungen. 1835}{magnetische_beobachtungen_1835}
\begin{center}
\hrule\vspace{2mm}
Göttingische gelehrte Anzeigen. 7. März 1835.
\vspace{2mm}\hrule\vspace{5mm}
\end{center}
In der Sitzung der Königlichen Sozietät am 14. Februar berichtete Hofrat \textsc{Gauss} über die Beobachtungen im magnetischen Observatorium und über die damit verbundenen Beobachtungen an anderen Orten. Wir teilen hier einen Auszug daraus mit, der als
\begin{center}\itshape Fortsetzung der am 9. August 1834 gegebenen Nachricht\end{center}
\noindent betrachtet werden kann.
\end{samepage}

\Needspace{13\baselineskip}
\SourceLeaf{535}Die zweimal täglich vorgenommenen Aufzeichnungen der Nadelstellung wurden ohne Unterbrechung fortgesetzt und umfassen jetzt bereits beinahe ein volles Jahr. Die Monatsmittel seit Juli des vorigen Jahres waren:
\begin{center}
\begin{tabular}{rl|rrr|rrr}
 & & \multicolumn{3}{c|}{8 Uhr vormittags} & \multicolumn{3}{c}{1 Uhr nachmittags}\\\hline
1834 & August & $18^0$ & $38'$ & $48''1$ & $18^0$ & $49'$ & $11''0$\\
 & September & & $36$ & $58.4$ & & $46$ & $32.3$\\
 & Oktober & & $37$ & $18.4$ & & $44$ & $47.2$\\
 & November & & $37$ & $38.4$ & & $43$ & $4.3$\\
 & Dezember & & $37$ & $54.8$ & & $41$ & $32.7$\\
1835 & Januar & & $37$ & $51.5$ & & $42$ & $14.4$
\end{tabular}
\end{center}

Die vereinbarten Beobachtungen an bestimmten Tagen in kurzen, ununterbrochenen Zeitabständen haben seit der letzten Nachricht an vier Hauptterminen stattgefunden, einige außerordentliche Nebentermine nicht mitgerechnet. In den letzten Monaten wurden an der Einrichtung dieser Beobachtungen einige Änderungen vorgenommen und an anderer Stelle bekannt gemacht. Die Teilnahme hat sich bereits ausgeweitet und wird sich bald noch weiter ausweiten. Die Beobachtungen haben auch schon sehr bemerkenswerte Ergebnisse geliefert, ähnlich denen, die im früheren Bericht erwähnt wurden. Eine grafische Darstellung der Übereinstimmung zwischen den Beobachtungen vom 1. und 2. Oktober und vom 29. und 30. November in Göttingen, Leipzig und Berlin wird demnächst in \textsc{Poggendorffs} Annalen der Physik erscheinen. Noch bemerkenswerter ist jedoch die Übereinstimmung der Beobachtungen vom 5. und 6. November in Kopenhagen und Mailand in allen zahlreichen und auffallend großen Schwankungen, von denen ebenfalls an anderer Stelle eine Zeichnung gegeben werden wird. Wir treten hier in eine Welt geheimnisvoller Naturkräfte ein, deren wunderbar wechselndes Spiel sich über die halbe Breite Europas im selben Augenblick und bis in die kleinsten Nuancen auf dieselbe Weise offenbart und für deren Wirkungskreis diese Standlinie noch viel zu kurz erscheint.

Die hiesigen Einrichtungen für magnetische Beobachtungen wurden inzwischen in mehreren wesentlichen Punkten erweitert. Für manche Beobachtungen ist, wenn hohe Genauigkeit gefordert wird, ein zweiter Apparat in einiger Entfernung vom Hauptapparat unerlässlich, damit man die stündlichen Veränderungen der magnetischen Kraft berücksichtigen kann. Zu diesem Zweck hängt seit August des vorigen Jahres ein großer Magnetstab in der Sternwarte, nachdem die im Jahr 1832 ver\SourceLeaf{536}wendeten Apparate an das physikalische Kabinett abgegeben worden waren. Sein übriges Zubehör ist dem des Stabes im magnetischen Observatorium ganz ähnlich. Der Magnetstab in der Sternwarte besteht ebenfalls aus Uslarschem Gussstahl, ist 4 Fuß lang, fast drei Zoll breit und über einen halben Zoll dick und wiegt 25 Pfund. Er hängt an einem 16 Fuß langen Seidenfaden aus tausend Einzelfäden\HistoricalNote{536}, der oberhalb der Saaldecke befestigt ist und frei durch eine kleine Öffnung läuft, die in der Decke angebracht wurde. Der unmittelbare Grund für die Wahl eines so schweren Stabes war die Absicht, den Luftzug unschädlich zu machen. Er lässt sich in diesem Raum nicht immer ganz fernhalten und hat die kleineren Apparate trotz des sie umschließenden Schutzkastens öfter gestört. Das Ergebnis entsprach \textit{dieser} Erwartung und übertraf die übrigen Erwartungen an die Genauigkeit aller damit möglichen Beobachtungen noch bei Weitem. Nur absolute Beobachtungen der Deklination und Intensität bleiben wegen des an vielen Stellen in der Sternwarte vorhandenen Eisens natürlich ausgeschlossen.

Die größte Schwingung, die der den Stab umschließende Kasten zulässt, beträgt etwa 27 Grad. Die größte Schwingung, die sich noch unmittelbar auf der Skala messen lässt, beträgt 9 bis 10 Grad, weil bei größeren Schwingungen die Visierlinie des Fernrohrs den fast vier Zoll breiten Spiegel nicht mehr trifft. Ist der Stab einmal in Schwingung versetzt, so nehmen die Schwingungen in geometrischer Folge so langsam ab, dass sie oft erst nach 10 oder mehr Stunden auf die Hälfte zurückgehen. Zuweilen geschieht das aber auch viel früher, worauf wir unten noch eigens eingehen werden. Der jetzt eingehängte Stab ist der stärkste aus einer größeren Anzahl von Stäben, die für das physikalische Kabinett angefertigt wurden. Eine Schwingung dauert etwas über 42 Sekunden. Diese Dauer unterliegt wegen der Temperatur und der Veränderlichkeit des Erdmagnetismus einigen, wenn auch sehr kleinen Veränderungen (und vielleicht kann sich auch die Kraft des Stabes selbst im Laufe der Zeit ändern, obwohl eine solche Änderung bisher noch gar nicht bemerkbar ist). Sie lässt sich schon aus wenigen Schwingungen so genau bestimmen, dass man den Stab dann 8 oder mehr Stunden sich selbst überlassen kann, ohne hinterher über die Anzahl der inzwischen vollendeten Schwingungen im Zweifel zu sein.

\SourceLeaf{537}Ebenso interessant wie die rein magnetischen Beobachtungen sind die elektrodynamischen Versuche, die sich mit diesem Apparat anstellen lassen. Zu diesem Zweck ist der Stab von einem ähnlichen Multiplikator umgeben wie der Stab im magnetischen Observatorium. Dieser Multiplikator ist jedoch größer und hat eine Drahtlänge von 2700 Fuß in 270 Windungen. Er ist in die große, schon im früheren Bericht erwähnte Drahtkette eingefügt, die die Sternwarte, das magnetische Observatorium und das physikalische Kabinett verbindet. Der galvanische Strom muss darin insgesamt 11000 Fuß Draht, also fast eine halbe geografische Meile, durchlaufen und wirkt dann zugleich auf drei magnetische Apparate, nämlich
\begin{list}{}{\setlength{\leftmargin}{2em}\setlength{\labelwidth}{2em}\setlength{\itemsep}{0pt}\setlength{\parsep}{0pt}}
\item[I.] den 25 Pfund schweren Stab in der Sternwarte,
\item[II.] den 4 Pfund schweren Stab im magnetischen Observatorium\newline (Multiplikator mit 200 Windungen)
\item[III.] den ein Pfund schweren Stab im physikalischen Kabinett\newline (Multiplikator mit 160 Windungen).
\end{list}
\noindent Einzelne Teile der Kette lassen sich nach Belieben in vielen Kombinationen leicht abtrennen.

Von den zahlreichen Versuchen, die mit diesen Apparaten bereits angestellt worden sind, führen wir hier nur einige an.

Wird ein galvanischer Strom mit der Kette verbunden, setzen sich die Magnetstäbe in den drei Apparaten so unmittelbar in Bewegung, dass sich der Beginn ihrer Bewegung bis auf einen kleinen Bruchteil einer Zeitsekunde genau beobachten lässt. Der Vergleich der Uhren bei den drei Apparaten liefert so vollkommen übereinstimmende Ergebnisse, ob der Strom am einen Ende, am anderen oder in der Mitte erzeugt wird, dass dies die Unmessbarkeit der Zeit, in der der Strom eine halbe Meile durchläuft, vollkommen bestätigt. Nach den interessanten Versuchen von \textsc{Wheatstone}, die kürzlich in den \textit{Philosophical Transactions} für 1834 veröffentlicht wurden und nach denen sich der elektrische Strom im Metall anscheinend schneller bewegt als das Licht im Raum, ließ sich ein solches Ergebnis freilich schon vermuten. Auf das Verhalten eines \textit{galvanischen} Stroms und seine Einwirkung auf die Magnetnadel konnte man daraus jedoch nur mit Vorbehalt schließen.

Die Intensität eines galvanischen Stroms wird durch die Ablenkung der Magnetnadel, also zunächst in Skalenteilen, gemessen oder bestimmt, allerdings
\SourceLeaf{538}offenbar in den drei Apparaten mit verschiedenen Einheiten, die von den Abmessungen der Multiplikatoren und dem Wert der Skalenteile in Bogensekunden abhängen. Zahlreiche Versuche zeigen nun, dass zwischen den Ablenkungen, die derselbe Strom im selben Augenblick an den drei Apparaten hervorruft, stets ein genau konstantes Verhältnis besteht, ob der Strom am einen Ende, am anderen oder in der Mitte erzeugt wird. Daraus ergibt sich das wichtige Ergebnis, dass der Strom auf seiner ganzen Länge dieselbe Intensität hat, jedenfalls nichts Merkliches davon verliert. Künftig wird man besonders darauf achten, ob dieses Ergebnis auch unter besonderen Umständen, vor allem bei starkem Regen, gültig bleibt.

Bei allen drei Apparaten sind Kommutatoren (Gyrotrope) mit der Kette verbunden, mit denen sich die Stromrichtung leicht umkehren lässt. Hofrat \textsc{Gauss} hat dem Kommutator in der Sternwarte eine besondere Einrichtung gegeben, sodass ein einziger Fingerdruck den Strom ganz augenblicklich umkehrt. Wiederholt man diese Umkehrung stets in Zeitabständen, die der Dauer einer Schwingung des einen Stabes entsprechen, so werden die Schwingungen dieses Stabes immer größer. Das wurde für einen Versuch genutzt, der eine auffallende mechanische Wirkung hervorbringt. Professor \textsc{Weber} ließ neben dem Magnetstab im physikalischen Kabinett einen leichten Auslösemechanismus für einen Wecker oder eine Pendeluhr anbringen. Der von der Sternwarte aus geleitete Strom betätigt sie nach ein paar Schwingungen jedes Mal vollkommen zuverlässig. Mit dem 25 Pfund schweren Stab ließe sich offensichtlich eine noch viel stärkere mechanische Wirkung erzielen.

Besonders wichtige Dienste leisten diese Apparate bei der Untersuchung der mathematischen Gesetze, nach denen sich die Erzeugung und Wirkung der von \textsc{Faraday} entdeckten magnetoelektrischen Induktion richten, und bei ihrer Zurückführung auf absolute Maße. Hofrat \textsc{Gauss} wird die Ergebnisse seiner Untersuchungen darüber zu gegebener Zeit an anderer Stelle bekannt machen. Von den dabei verwendeten Vorrichtungen erwähnen wir hier nur eine, mit der sich diese Induktion ebenso einfach wie genau messbar hervorrufen lässt. Um eine hölzerne Spule ist ein umsponnener Draht in 1050 Windungen geführt, dessen Enden durch den Kommutator mit der Kette verbunden wer\SourceLeaf{539}den. Die Spule lässt sich über die freistehende Hälfte eines starken Magnetstabes schieben. Dabei fließt jedes Mal ein galvanischer Strom durch die Kette, ein starker von kurzer Dauer oder ein schwächerer von längerer Dauer, je nachdem, ob die Bewegung schneller oder langsamer erfolgt, sodass die Gesamtwirkung eines einzigen Aufschiebens von der Geschwindigkeit der Bewegung unabhängig ist. Der Strom selbst dauert immer nur so lange wie die Bewegung der Spule. Das Abziehen der Spule erzeugt einen entgegengesetzten Strom, ebenso das Aufschieben mit dem entgegengesetzten Ende. Geschieht die Bewegung sehr schnell, so ist die Wirkung des Stroms auf die Magnetnadel in einem der mit der Kette verbundenen Multiplikatoren einem augenblicklichen Stoß von bestimmter Stärke gleichzusetzen. Abziehen und umgekehrtes Wiederaufstecken bewirken also zwei gleichgerichtete Impulse der Magnetnadel. Erneutes Abziehen und wieder umgekehrtes Aufschieben würde daher zwei untereinander gleiche, den vorherigen aber entgegengesetzte Impulse hervorbringen. Wird dazwischen der Kommutator umgestellt, so erfolgen auch die letzten beiden Wirkungen im selben Sinn wie die ersten beiden. Ein solcher vollständiger Wechsel (Abziehen, umgekehrtes Aufstecken und Umstellen des Kommutators) gelingt bequem in zwei Sekunden. Wenn man will, kann man deshalb während einer Schwingung des großen Magnetstabes bequem und im Takt 21 Wechsel ausführen und den Stab dadurch so stark in Bewegung setzen, dass die ganze Skala aus dem Gesichtsfeld des Fernrohrs gerät. Dieser Hinweis wird genügen, um zu erkennen, wie sich die Stärke des durch diese Art der Induktion entstehenden galvanischen Stroms genau messen lässt. Diese Stärke hängt aber zugleich vom Widerstand der Kette selbst ab und nimmt mehr oder weniger zu, je nachdem, ob mehr oder weniger Stücke der Kette abgetrennt werden. Auf diese Weise wurde das Verhältnis der Widerstände in den einzelnen Bestandteilen der Kette und den Multiplikatoren sehr genau bestimmt. Durch vielfältige Kombinationen wurde das schöne, von \textsc{Ohm} aufgestellte Gesetz, dem die Intensität eines Stroms bei einer Teilung folgt, aufs Vollkommenste bestätigt. Nahezu übereinstimmende Ergebnisse wurden auch mit hydrogalvanischen Strömen gefunden. Diese eignen sich allerdings wegen ihrer veränderlichen Stärke weniger für solche Bestimmungen und erfordern bei den Versuchen deshalb jedenfalls noch besondere Vorsichtsmaßnahmen. Vielleicht ist es nicht ohne Interesse, wenn hier bemerkt
\SourceLeaf{540}wird, dass der gesamte Widerstand in der durch die Luft geführten doppelten Drahtverbindung zwischen der Sternwarte und dem physikalischen Kabinett bei einer Drahtlänge von mehr als 6000 Fuß nur etwa halb so groß ist wie der Widerstand, auf den der Strom allein im Multiplikator des magnetischen Observatoriums (Drahtlänge 1100 Fuß) trifft, oder nur ein Sechstel des Widerstands der gesamten Kette beträgt. Das erklärt sich leicht aus der unterschiedlichen Dicke des Drahts. Alle Versuche bestätigen, dass bei Drähten aus demselben Metall der Widerstand immer direkt proportional zur Länge und umgekehrt proportional zur \textit{Fläche} des Querschnitts ist.

Wir haben oben erwähnt, dass die Abnahme des Schwingungsbogens der großen Nadel zu verschiedenen Zeiten sehr unterschiedlich war. Ähnliche Unterschiede hatten sich bereits 1832 an den kleinen Apparaten und später auch an der Nadel im magnetischen Observatorium gezeigt. Diese Unterschiede bleiben aber stets in viel engeren Grenzen als beim Stab der Sternwarte, wo die Abnahme des Schwingungsbogens von einer Schwingung zur nächsten in verschiedenen Versuchsreihen zwischen $\frac{1}{5000}$ und $\frac{1}{50}$ schwankte. Diese bemerkenswerte Erscheinung hat die Aufmerksamkeit von Hofrat \textsc{Gauss} besonders auf sich gezogen. Es scheinen dabei \textit{mehrere} Ursachen zusammenzuwirken, die zum Teil noch heute rätselhaft bleiben. Inzwischen ist es Hofrat \textsc{Gauss} gelungen, die Ursache mit dem weitaus stärksten Einfluss zu ermitteln. Er bemerkte nämlich, dass der Schwingungsbogen bei geschlossener Kette stets viel schneller abnahm als bei offener. So war es leicht, als Ursache dieser schnellen Abnahme die Rückwirkung eines galvanischen Stroms zu erkennen, der durch die Schwingung der Nadel selbst mittels Induktion in der Kette erzeugt wird. Bei der folgenden Rückschwingung hat der Strom die entgegengesetzte Richtung und wirkt stets auf eine Verringerung des Schwingungsbogens hin. Diese Erklärung bestätigte sich vollkommen. Der Schwingungsbogen nahm bei offener Kette am langsamsten ab, bei geschlossener, aber vollständiger Kette schneller, bei abgetrennten einzelnen Stücken der Kette noch schneller und am schnellsten (sodass er in einer halben Stunde auf die Hälfte zurückging), wenn die Kette unmittelbar hinter dem Multiplikator des großen Stabes geschlossen war. Diese Unterschiede richteten sich sogar ganz nach der Größe des wirksam gebliebenen Teils der Kette.

\SourceLeaf{541}Nachdem diese Erklärung gefunden war, ließ sich das Ergebnis einiger Versuche leicht voraussehen, die wohl zu den auffallendsten im Bereich des Elektromagnetismus gerechnet werden dürfen. Selbst die quantitativen Verhältnisse der Erscheinungen ließen sich im Voraus berechnen und wurden bei den wiederholt angestellten Versuchen stets aufs Vollkommenste bestätigt. Es sind die folgenden.

Wird der Magnetstab in der Sternwarte (I) in Schwingungen versetzt, etwa so große, wie der Kasten zulässt, so haben diese gar keinen Einfluss auf die Nadeln im magnetischen Observatorium (II) oder im physikalischen Kabinett (III). Diese \textit{bleiben} in Ruhe, wenn sie vorher in Ruhe waren, vorausgesetzt, dass die Kette offen ist oder wenigstens die Multiplikatoren, die die beiden zuletzt genannten Nadeln umgeben, davon abgetrennt sind. In dem Augenblick aber, in dem die Kette geschlossen oder zum Beispiel der Multiplikator von II in die geschlossene Kette eingefügt wird, beginnt die Nadel II sofort mitzuschwingen. Hat die Nadel II schon vorher geschwungen, so erhalten ihre Schwingungen den besonderen Charakter \textit{gemischter} Schwingungen. Die eine hängt vom Anfangszustand ab und hat dieselbe Periode wie die Schwingungen dieser Nadel unter dem alleinigen Einfluss des Erdmagnetismus ($20''$). Die andere folgt einer Periode von $42''$ (wie die große Nadel I) und ist ihrer Größe nach dem Schwingungsbogen von I proportional (etwa $\frac{1}{200}$, wenn die Kette hinter dem Multiplikator von II abgetrennt ist). Das stimmt vollkommen mit den Ergebnissen der Theorie überein, ebenso wie der stets genau bestätigte Umstand, dass die Schwingungen von I und die induzierten Schwingungen von II bei gleich langen Perioden einen verschiedenen Anfang haben. Sie unterscheiden sich darin stets um eine halbe Schwingungszeit ($21''$), und zwar in dem Sinn, den die Theorie für die gegebenen Umstände vorhersagt. Was hier am Beispiel der Nadel II gesagt wurde, gilt ganz ähnlich für die Nadel III. Ihre natürliche Schwingungsdauer beträgt $14''$, und unter der Wirkung der Induktion führt sie zusammengesetzte Schwingungen mit Perioden von $14''$ und $42''$ aus.

Ein ganz anderes Ergebnis muss nach der Theorie eintreten, wenn sich eine zweite Nadel mit einem Multiplikator in der Kette befindet, in der der große Stab schwingt, und ihre natürliche Schwingungsdauer genau so groß ist wie die des großen Magnetstabes. Diese zweite Nadel, die völlig in Ruhe ist,
\SourceLeaf{542}solange die Kette offen ist, beginnt ebenfalls in dem Augenblick mitzuschwingen, in dem die Kette geschlossen wird. Diese Schwingungen von derselben Dauer wie die natürlichen nehmen aber beständig an Größe zu, bis sie (erst nach sehr langer Zeit) ein Maximum erreichen, bei dem der Luftwiderstand die Zunahme durch die Induktionskraft ausgleicht. Um diesen bemerkenswerten Versuch tatsächlich durchführen zu können, wurde (da sich zu diesem Zeitpunkt mangels eines zweiten passenden Multiplikators kein großer Stab aufhängen ließ) der ein Pfund schwere Stab des physikalischen Kabinetts auf bekannte Weise durch Verbindung mit einem ähnlichen, etwas schwächer magnetisierten Stab astatisch gemacht, oder genauer gesagt, zu einer Doppelnadel, deren natürliche Schwingungsdauer genau auf $42''3$ gebracht wurde. Der Versuch gelang damit aufs Vollkommenste. Der in der Sternwarte schwingende Stab teilte der Doppelnadel im physikalischen Kabinett in dem Augenblick, in dem die Kette geschlossen wurde, wie durch eine wunderbare Sympathie seine Schwingungen mit, und zwar so, dass jede folgende um etwa 50 Skalenteile oder einen halben Grad größer wurde als die vorhergehende. Bald verschwand das gesamte Skalenbild aus dem Gesichtsfeld. An der immer größeren Geschwindigkeit, mit der es das Gesichtsfeld durchlief, konnte man aber weiterhin die Zunahme des Schwingungsbogens erkennen. Über eine Stunde lang wurde dieses wunderbare Spiel der Sympathie beobachtet.

Es braucht kaum bemerkt zu werden, dass auch der vier Pfund schwere Stab im magnetischen Observatorium in der geschlossenen Kette einen Strom induziert, dessen Vorhandensein an der schnellen Abnahme des Schwingungsbogens ganz eindeutig zu erkennen ist. Er muss daher auch auf die beiden anderen Stäbe ähnlich wirken, wie es der erste Versuch gezeigt hat. Die Rechnung ergibt jedoch, und die Erfahrung bestätigt, dass diese Wirkungen zu klein sind, um wahrnehmbar zu sein. Noch weniger könnte also der schwächste der drei Stäbe wahrnehmbare Wirkungen dieser Art erzeugen.

\clearpage
